% WACV 2026 Supplementary Material Template
% based on the WACV 2026 main paper template

\documentclass[10pt,twocolumn,letterpaper]{article}

%%%%%%%%% PAPER TYPE - PLEASE UPDATE FOR FINAL VERSION
%\usepackage[review,algorithms]{wacv}      % To produce the REVIEW version for the algorithms track
\usepackage[review,applications]{wacv}      % To produce the REVIEW version for the applications track
%\usepackage{wacv}              % To produce the CAMERA-READY version
%\usepackage[pagenumbers]{wacv} % To force page numbers, e.g. for an arXiv version

% It is strongly recommended to use hyperref, especially for the review version.
% hyperref with option pagebackref eases the reviewers' job.
% Please disable hyperref *only* if you encounter grave issues, 
% e.g. with the file validation for the camera-ready version.
%
% If you comment hyperref and then uncomment it, you should delete *.aux before re-running LaTeX.
% (Or just hit 'q' on the first LaTeX run, let it finish, and you should be clear).
\definecolor{wacvblue}{rgb}{0.21,0.49,0.74}
\usepackage[pagebackref,breaklinks,colorlinks,allcolors=wacvblue]{hyperref}

%%%%%%%%% PAPER ID - PLEASE UPDATE
\def\wacvPaperID{967} % *** Enter the WACV Paper ID here
\def\confName{WACV}
\def\confYear{2026}

%%%%%%%%% TITLE - PLEASE UPDATE
\title{Manga109Script Dataset for Script-to-Layout Generation}

%%%%%%%%% AUTHORS - PLEASE UPDATE
\author{First Author\\
Institution1\\
Institution1 address\\
{\tt\small firstauthor@i1.org}
\and
Second Author\\
Institution2\\
First line of institution2 address\\
{\tt\small secondauthor@i2.org}
}

\begin{document}

%%%%%%%%% SUPPLEMENTARY MATERIAL TITLE
\maketitlesupplementary
\renewcommand{\thesection}{\Alph{section}}
\renewcommand{\thesubsection}{\thesection.\arabic{subsection}}

\section{Detailed Implementation Parameters}

As referenced in Section 4.4 of the main paper, we provide comprehensive hyperparameters for reproducing our Script2Layout models. All models were fine-tuned using LoRA with the Hugging Face \texttt{transformers} and \texttt{peft} libraries on a single NVIDIA H100 GPU.

\subsection{Dataset Configuration}

\begin{itemize}
    \item \textbf{Maximum token length:} 4096
\end{itemize}

\subsection{Model Configuration}

\begin{itemize}
    \item \textbf{Data type:} \texttt{torch.float16}
    \item \textbf{Attention implementation:} FlashAttention 2
    \item \textbf{Tokenizer settings:} Right padding, EOS token addition enabled
\end{itemize}

\subsection{LoRA Configuration}

\begin{itemize}
    \item \textbf{Reduction factor:} $r = 8$
    \item \textbf{LoRA alpha:} 16
    \item \textbf{Target modules:} \texttt{q\_proj} and \texttt{v\_proj}
    \item \textbf{LoRA dropout:} 0.1
    \item \textbf{Bias setting:} \texttt{none}
    \item \textbf{Task type:} \texttt{CAUSAL\_LM}
\end{itemize}

\subsection{Training Parameters}

\begin{itemize}
    \item \textbf{Batch size:} 1 per device (training and evaluation)
    \item \textbf{Gradient accumulation steps:} 32
    \item \textbf{Training epochs:} 20
    \item \textbf{Optimizer:} \texttt{adamw\_torch\_fused}
    \item \textbf{Learning rate:} $2 \times 10^{-4}$
    \item \textbf{LR scheduler:} Cosine scheduler
    \item \textbf{Mixed precision:} FP16 enabled
    \item \textbf{Logging frequency:} Every 100 steps
    \item \textbf{Evaluation strategy:} Every 100 steps with data grouped by length
    \item \textbf{Save strategy:} Every epoch
    \item \textbf{Experiment tracking:} Weights \& Biases (W\&B)
\end{itemize}

\end{document}